% !TeX program = xelatex
% Научный отчёт о текущем состоянии проекта voice_signals
% Дата: 23 октября 2026
% Назначение: материал для обсуждения с научным руководителем
\documentclass[12pt,a4paper]{article}

\usepackage{fontspec}
\setmainfont{XITS}
\setsansfont{DejaVu Sans}
\setmonofont{DejaVu Sans Mono}
\usepackage{polyglossia}
\setdefaultlanguage{russian}
\setotherlanguage{english}
\usepackage[a4paper,margin=2cm]{geometry}
\usepackage{amsmath,amssymb}
\usepackage{graphicx}
\usepackage{booktabs}
\usepackage{longtable}
\usepackage{array}
\usepackage{hyperref}
\usepackage{enumitem}
\usepackage{xcolor}
\usepackage{listings}
\usepackage{caption}

\hypersetup{
  colorlinks=true,
  linkcolor=blue,
  urlcolor=blue,
  citecolor=blue,
  pdftitle={Отчёт о состоянии проекта voice\_signals},
  pdfauthor={}
}

\lstset{
  basicstyle=\ttfamily\small,
  breaklines=true,
  frame=single,
  columns=fullflexible,
  backgroundcolor=\color{gray!5}
}

% Снижаем overfull hbox от длинных неразбиваемых идентификаторов кода
% (snake_case) внутри inline \texttt{...}, допуская больше растяжения строк.
\sloppy
\emergencystretch=3em
\hbadness=10000

\title{\textbf{Анализ временной динамики речевых сигналов \\
для выявления интеллектуальных способностей человека: \\
отчёт о текущем состоянии проекта \texttt{voice\_signals}}}
\author{}
\date{23 октября 2026 г.}

\begin{document}

\maketitle

\begin{abstract}
\noindent
Отчёт описывает текущее состояние реализации проекта \texttt{voice\_signals} —
исследовательского программного инструмента для детекции речевой активности
(VAD) и анализа нелинейной временной динамики речевого сигнала методом
модифицированного нормированного размаха (modified R/S, методика
Кузенкова и Логинова, 2014). Представлены: цель работы и научная гипотеза,
обзор методической базы, архитектура и состав реализованного программного
обеспечения, результаты пилотных вычислительных экспериментов на синтетических
данных и на одной тестовой записи, обнаруженные методические ограничения
(зависимость оценок показателя Хёрста $H$ и характерного времени $T$ от длины
анализируемого сегмента и от окна огибающей), а также перечень критических
блокеров, препятствующих переходу к содержательной научной оценке.
Документ подготовлен как основа для обсуждения с научным руководителем:
зафиксировано, что сделано, что работает по инженерным критериям, что
проверено лишь технически, и какие решения нужны для продолжения работы.
\end{abstract}

\tableofcontents

% =====================================================================
\section{Введение и постановка задачи}
% =====================================================================

\subsection{Научная идея}

Работа развивает гипотезу о том, что временная структура речевого сигнала
человека (чередование речевых и паузальных участков, их статистические
характеристики) связана с его когнитивными/интеллектуальными способностями.
Методической основой служит модифицированный метод нормированного размаха
(R/S-анализ), предложенный в статье Н.\,П.~Кузенкова и В.\,М.~Логинова
\cite{kuzenkov2014}, где показатель Хёрста $H$ и характерное время $T$
вычисляются по огибающей мощности речевого сигнала и используются для
различения речевых патологий неврологического генеза (афазии, дизартрии)
от нормы.

Проект \texttt{voice\_signals} является инструментальным слоем для этой
гипотезы. Он решает три последовательные инженерные подзадачи:

\begin{enumerate}
  \item \textbf{Сегментация речи} (VAD) — отделение речевых участков от пауз,
  хезитаций и артефактов;
  \item \textbf{Расчёт нелинейных динамических признаков} ($H$, $T$,
  в перспективе DFA-оценка Хёрста, показатель Ляпунова, энтропийные меры) на
  найденных речевых сегментах;
  \item \textbf{Проверка связи признаков с результатами когнитивных тестов}
  — отложена до накопления реального корпуса с разметкой.
\end{enumerate}

\subsection{Принципиальное разграничение (методологическая позиция)}

В правилах проекта зафиксировано явное разделение технической и содержательной
задачи:

\begin{quote}
\itshape
Модель сегментирует речевые события. Полученные признаки могут использоваться
для дальнейшего исследования ассоциаций с результатами когнитивных тестов.
\end{quote}

То есть на текущем этапе VAD и расчёт $H/T$ \textbf{не интерпретируются} как
модель интеллекта или когнитивного состояния — это сознательное ограничение,
принятое для методологической чистоты. Любые клинические или когнитивные
выводы статьи~\cite{kuzenkov2014} не переносятся на проект автоматически.

\subsection{Текущий статус}

На момент подготовки отчёта (23.10.2026) завершены две итерации разработки
(v0.1.0 и v0.2.0). Инженерная база функциональна и покрыта тестами, но
\textbf{отсутствует доступ к реальному речевому корпусу с разметкой}, что
является главным блокером перехода от инженерной проверки работоспособности
к научной оценке качества и к проверке целевой гипотезы.

% =====================================================================
\section{Методическая основа: модифицированный R/S-анализ}
\label{sec:method}
% =====================================================================

\subsection{Классический R/S-анализ}

Для временного ряда $x(t_i)$ на интервале длины $\tau$ классический
R/S-анализ использует:

\begin{align}
\bar{x}(\tau) &= \frac{1}{\tau}\sum_{i=1}^{\tau} x(t_i)
  && \text{(среднее)} \\
S(\tau) &= \sqrt{\frac{1}{\tau}\sum_{i=1}^{\tau}\big(x(t_i)-\bar{x}(\tau)\big)^2}
  && \text{(стандартное отклонение)} \\
X(t,\tau) &= \sum_{j=1}^{t}\big(x(j)-\bar{x}(\tau)\big)
  && \text{(накопившееся отклонение)} \\
R(\tau) &= \max_{1\le t\le \tau} X(t,\tau) - \min_{1\le t \le \tau} X(t,\tau)
  && \text{(размах)}
\end{align}

Эмпирический закон Хёрста:
\begin{equation}
\frac{R(\tau)}{S(\tau)} = \left(\frac{\tau}{2}\right)^{H},
\label{eq:hurst-law}
\end{equation}
где $H$ — показатель Хёрста. Интерпретация: $H=0.5$ соответствует
броуновскому движению (отсутствие памяти), $0.5<H<1.0$ — персистентному
процессу (тренды сохраняются), $0<H<0.5$ — антиперсистентному (возврат
к среднему).

Классическая схема имеет три проблемы: экспоненциальный рост плотности
отсчётов при равномерном переборе $\tau$ в логарифмической шкале,
чувствительность к начальной фазе сигнала и вычислительную сложность
$O(N^2)$.

\subsection{Модификация Кузенкова--Логинова}

В \cite{kuzenkov2014} предложена модификация:

\begin{itemize}
  \item шкала $\tau$ строится как геометрическая прогрессия
    $\tau_1=3,\ \tau_i = a\,\tau_{i-1}$, $a=(N/3)^{1/(n-1)}$ — равномерный шаг
    в логарифмических координатах;
  \item ряд разбивается на $k$ пересекающихся интервалов, $R/S$ усредняется по
    ним — снижается зависимость от фазы;
  \item вычислительная сложность снижается до $O(n\cdot N)$;
  \item в качестве исходного ряда берётся \emph{огибающая} речевого сигнала —
    среднеквадратичное отклонение в скользящем окне (в оригинальной работе
    100~мс при частоте дискретизации 8~кГц);
  \item зависимость $\ln(R/S)$ от $\ln(\tau/2)$ аппроксимируется
    кусочно-полиномиальной функцией для оценки характерного времени $T$ —
    момента смены режима (перехода от одного наклона к другому).
\end{itemize}

В клиническом исследовании~\cite{kuzenkov2014} (429 образцов, суммарно
23\,ч~56\,мин записи, рассказ А.~П.~Чехова «Злой мальчик», чтение и
пересказ) получены оценки, сведённые в таблицу~\ref{tab:kuzenkov-results}.

\begin{table}[h]
\centering
\caption{Статистические характеристики $H$ и $T$ по группам
(из \cite{kuzenkov2014})}
\label{tab:kuzenkov-results}
\small
\begin{tabular}{lcccccc}
\toprule
Группа & Корр. & $H$ (ср.) & $H$ (СКО) & $T$ (ср., с) & $T$ (СКО) \\
\midrule
Афазия, чтение        & $-0.643$ & 0.989 & 0.00488 & 1.21  & 0.403 \\
Афазия, пересказ       & $-0.610$ & 0.990 & 0.00509 & 1.06  & 0.642 \\
Дизартрия, чтение      & $-0.686$ & 0.988 & 0.00511 & 1.33  & 0.532 \\
Дизартрия, пересказ    & $-0.405$ & 0.988 & 0.00488 & 0.966 & 0.638 \\
Норма, чтение          & $-0.850$ & 0.990 & 0.00289 & 0.902 & 0.259 \\
Норма, пересказ        & $-0.852$ & 0.991 & 0.00384 & 0.934 & 0.497 \\
Нейродин. афазия, чтение    & $-0.289$ & 0.985 & 0.00593 & 1.28 & 0.467 \\
Нейродин. афазия, пересказ  & $-0.391$ & 0.984 & 0.00671 & 0.997 & 0.576 \\
\bottomrule
\end{tabular}
\end{table}

Статья показывает приблизительно линейную связь
$T \approx \alpha\,(1-H)$ (уравнения регрессий $T=-kH+b$ для каждой группы) и
использует рандомизационный тест (5000 итераций) для проверки значимости
различий с нормой: достоверные различия ($p<0.05$) выявлены для дизартрии
и нейродинамической афазии, тогда как моторная афазия в режиме монолога
статистически близка к норме.

\textbf{Важно для дальнейшего обсуждения:} во всех группах $H$ лежит в
узком диапазоне $0.984$--$0.991$, то есть основная различающая переменная —
$T$, а не $H$. Это наблюдение существенно для интерпретации собственных
результатов проекта (раздел~\ref{sec:experiment}).

% =====================================================================
\section{Архитектура и реализация}
\label{sec:architecture}
% =====================================================================

Проект реализован как локальный Python-пакет \texttt{voice\_signals}
(версия 0.2.0, вторая итерация завершена 17.10.2026) с минимальными
зависимостями: единственная внешняя библиотека — \texttt{pyyaml};
NumPy, SciPy, PyTorch, librosa не используются. Все алгоритмы написаны
на чистом Python и детерминированы.

\subsection{Структура пакета}

\begin{lstlisting}
voice_signals/
  src/voice_signals/
    vad.py           # EnergyVad: энергетический VAD с гистерезисом
    vad_backend.py    # Унифицированный интерфейс VadBackend
    dynamics.py        # modified R/S: Hurst H, характерное время T
    framing.py         # Единая временная сетка, RMS-огибающая
    schemas.py          # Контракты данных (SegmentRecord и др.)
    config.py            # YAML-конфигурация, хэш конфигурации (SHA256)
  scripts/           # run_vad.py, run_modified_rs.py, run_audit.py,
                      # generate_synthetic.py, extract_features.py,
                      # demo_pipeline.py
  tests/              # 15 детерминированных unit-тестов (pytest)
  configs/default.yaml
  data/{synthetic,real,raw,segments,features}/
  docs/, reports/
\end{lstlisting}

\subsection{Единый слой фрейминга}

Класс \texttt{TimeGrid} (\texttt{framing.py}) задаёт общую временную сетку
(размер окна и шаг в отсчётах, вычисляемые из миллисекунд и частоты
дискретизации) и используется одновременно VAD-модулем и модулем расчёта
динамических признаков, что исключает рассинхронизацию временных меток между
подсистемами. Функция \texttt{compute\_rms\_frames} вычисляет RMS-огибающую
сигнала по этой сетке.

\subsection{VAD-модуль (Energy VAD)}

Базовый детектор речевой активности \texttt{EnergyVad} — энергетический VAD
с двухпороговым гистерезисом:

\begin{itemize}
  \item кадр 20~мс, шаг 10~мс;
  \item порог входа в речь: превышение энергии над оценкой фона (20-й
   перцентиль распределения энергии по файлу) на 10~дБ;
  \item порог выхода из речи: превышение на 6~дБ (гистерезис устраняет
   дребезг на границе речь/пауза);
  \item постобработка: отбрасывание речевых фрагментов короче 80~мс,
   склейка пауз короче 80~мс между соседними речевыми фрагментами.
\end{itemize}

Алгоритм детерминирован, не использует внешних ML-моделей, и выдаёт список
\emph{непересекающихся} сегментов \texttt{SegmentRecord} — унифицированного
контракта данных, валидируемого функцией
\texttt{validate\_non\_overlapping}. Класс наследуется от абстрактного
\texttt{VadBackend}, что подготавливает единый интерфейс для будущих
адаптеров WebRTC VAD и Silero VAD (на момент отчёта не реализованы — есть
только интерфейс).

\subsection{Модуль нелинейной динамики (modified R/S)}

Функция \texttt{modified\_rs\_hurst} (\texttt{dynamics.py}) реализует
алгоритмическую схему из раздела~\ref{sec:method} на энергетической
огибающей сигнала:

\begin{enumerate}
  \item построение огибающей как RMS в скользящем немедленно соседствующем
    окне (\texttt{envelope\_window\_sec}, по умолчанию 25~мс — отличается
    от 100~мс в оригинальной статье, см.\ обсуждение в
    разделе~\ref{sec:experiment});
  \item заявленная на момент написания отчёта «геометрическая шкала
    масштабов $\tau$» от \texttt{min\_tau} (по умолчанию 4) до половины
    длины огибающей \textbf{не соответствовала коду}: на практике
    использовался арифметический \texttt{range(...)}. Расхождение
    зафиксировано и устранено 31.10.2026, см.\ erratum в конце
    раздела~\ref{sec:experiment};
  \item для каждого $\tau$ — усреднение $R/S$ по перекрывающимся окнам
    (шаг $\tau/2$);
  \item оценка $H$ как наклон линейной регрессии $\ln(R/S)$ от
    $\ln(\tau/2)$ (метод наименьших квадратов, без внешних библиотек);
  \item оценка качества аппроксимации через $R^2$; при $R^2<0.7$
    устанавливается предупреждение \texttt{weak\_log\_log\_fit};
  \item оценка характерного времени $T$ двумя методами, управляемыми флагом
   \texttt{use\_piecewise\_fit} (по умолчанию \texttt{True}):
   \begin{itemize}
     \item \texttt{"argmax"} (итерация~01) — $\tau$, максимизирующий
       усреднённое значение $R/S$;
     \item \texttt{"piecewise"} (итерация~02, реализовано в
       \texttt{\_piecewise\_fit\_characteristic\_time}) — полный перебор
       точки излома, минимизирующей суммарный остаток двух независимых
       линейных МНК-регрессий (левого и правого сегмента на
       $\ln(\tau/2)$--$\ln(R/S)$); сохранена обратная совместимость.
   \end{itemize}
\end{enumerate}

Результат оформлен как неизменяемая структура \texttt{HurstRSResult} с
полями \texttt{hurst\_h}, \texttt{characteristic\_time\_sec},
\texttt{tau\_values}, \texttt{fit\_r\_squared}, \texttt{valid},
\texttt{warnings}, \texttt{fit\_method}. Короткие ($<16$ отсчётов) или
вырожденные (постоянные) входы дают \texttt{valid=False} с явным кодом
причины, а не ошибочное число.

\subsection{Контракты данных}

Пакет \texttt{schemas.py} фиксирует общие структуры для всех подсистем:
\texttt{SegmentRecord} (речевой сегмент с валидацией на непересечение и
допустимость метки), \texttt{FileAuditRecord} (результат аудита WAV),
\texttt{FrameFeatureRecord} (метаданные кадровых признаков),
\texttt{ExperimentRecord} (результат эксперимента с train/val/test-разбиением
по респондентам, метриками и хэшем конфигурации). Унификация контрактов
предназначена для того, чтобы baseline-детекторы (energy, WebRTC) и
будущий собственный VAD возвращали данные в одном формате и были сравнимы.

\subsection{Вспомогательный инструментарий}

\begin{itemize}
  \item \textbf{Генератор синтетического корпуса}
    (\texttt{scripts/generate\_synthetic.py}) — создаёт WAV-файлы с
    известной (ground truth) разметкой из четырёх типов событий
    (\texttt{speech}, \texttt{silent\_pause}, \texttt{hesitation},
    \texttt{artifact}), контролируемым отношением сигнал/шум (5/10/20/30~дБ)
    и известным истинным $H$ (\texttt{true\_h}) — используется для проверки
    корректности алгоритмов без доступа к реальным записям;
  \item \textbf{Аудит WAV-файлов} (\texttt{scripts/run\_audit.py}) —
    метаданные (частота дискретизации, число каналов, битность,
    длительность), энергетические характеристики (RMS, пик, crest factor,
    уровень шума) и флаги качества (клиппинг, низкая энергия, почти полная
    тишина, нестандартная частота дискретизации);
  \item \textbf{Экстрактор кадровых признаков}
    (\texttt{scripts/extract\_features.py}) — RMS, энергия в дБ,
    zero-crossing rate, упрощённая спектральная плоскостность (без
    FFT-реализации, через отношение дисперсии к максимуму — заявленное
    ограничение, не полноценный спектральный признак);
  \item \textbf{Демонстрационный конвейер} (\texttt{scripts/demo\_pipeline.py})
    — последовательный прогон: загрузка конфигурации $\to$ VAD $\to$
    modified R/S на найденном сегменте $\to$ интерпретация $H$.
\end{itemize}

\subsection{Воспроизводимость}

Конфигурация (\texttt{configs/default.yaml}) хэшируется функцией
\texttt{config\_hash} (SHA256) для привязки экспериментов к точному набору
параметров. Все алгоритмы детерминированы (повторный запуск даёт идентичный
результат — проверено unit-тестом
\texttt{test\_hurst\_result\_is\_deterministic}).

% =====================================================================
\section{Тестирование}
% =====================================================================

На момент отчёта в проекте 15 unit-тестов (\texttt{pytest}), организованных
в три модуля:

\begin{itemize}
  \item \texttt{test\_vad.py} — проверка энергетического VAD на синтетическом
   тоне (сегментация точно совпадает с границами тона с допуском 20~мс),
   отбрасывание слишком коротких речевых фрагментов, валидация контракта
   \texttt{SegmentRecord} (отказ при пересечении границ и недопустимой
   метке), детерминированность \texttt{modified\_rs\_hurst};
  \item \texttt{test\_framing.py} — корректность построения временной сетки,
   расчёта времён кадров, вычисления RMS-огибающей (включая проверку
   ожидаемого RMS синусоиды $A/\sqrt2$) и детерминированность;
  \item \texttt{test\_config.py} — загрузка YAML-конфигурации, детерминизм и
   чувствительность хэша конфигурации к изменению содержимого, валидация
   обязательных полей.
\end{itemize}

Все 15 тестов проходят (\texttt{15 passed}). Это подтверждает корректность
реализации алгоритмов \emph{относительно их собственной спецификации} на
контролируемых синтетических сигналах — не более того: тесты не могут
подтвердить качество VAD или релевантность признаков $H/T$ на реальной речи.

% =====================================================================
\section{Вычислительный эксперимент на тестовой записи}
\label{sec:experiment}
% =====================================================================

\subsection{Базовый прогон VAD}

На тестовой записи \texttt{tests/test\_data/Ru-Vegetables.wav} энергетический
VAD обнаружил 79 речевых сегментов длительностью от десятков миллисекунд до
нескольких секунд (полный список воспроизводимо сохранён в
\texttt{tests/result1.md} и совпадает между итерациями — подтверждение
детерминированности пайплайна). Для первого сегмента (1.29--1.94~с, то есть
длительностью 0.65~с) модифицированный R/S-анализ даёт:
\[
H = 0.979, \qquad T = 0.250\ \text{с}, \qquad R^2 = 0.995,
\qquad \text{fit\_method} = \text{"piecewise"}.
\]

Высокое значение $H$ ожидаемо и согласуется с \cite{kuzenkov2014}: там
$H$ во всех клинических группах тоже лежит в диапазоне 0.984--0.991
(таблица~\ref{tab:kuzenkov-results}), поскольку расчёт ведётся не по
сырому сигналу, а по гладкой огибающей мощности, которая по построению
обладает сильной положительной автокорреляцией.

\subsection{Обнаруженная методическая проблема: смещение оценок длиной
сегмента}

Был проведён дополнительный эксперимент: расчёт $H$ и $T$ на всех 79 найденных
речевых сегментах записи. Результаты:

\begin{itemize}
  \item короткие сегменты дают $H \approx 0.95$--$0.99$;
  \item длинные сегменты дают $H \approx 0.69$--$0.87$;
  \item корреляция $H$ с длиной сегмента $\approx -0.76$;
  \item корреляция $T$ с числом доступных масштабов $\tau$
   $\approx +0.87$ (то есть $T$ фактически воспроизводит длину сегмента, а
   не является от неё независимой характеристикой).
\end{itemize}

\textbf{Рабочая гипотеза причины.} Окно огибающей по умолчанию равно 25~мс
(в 4 раза короче, чем 100~мс в оригинальной статье), а верхняя граница
масштабов $\tau$ автоматически устанавливается как половина длины огибающей
сегмента — то есть напрямую зависит от длительности входного фрагмента.

Для проверки гипотезы окно огибающей было временно увеличено до 100~мс (как
в \cite{kuzenkov2014}) и прогон повторён. Результат:

\begin{itemize}
  \item значения $H$ снизились и стали более разбросанными ($\approx 0.4$--$0.9$);
  \item $T$ приблизилось к диапазону статьи ($\approx 0.6$--$2.0$~с);
  \item сегменты короче $\approx 1.2$~с перестали давать валидную оценку
   (отбрасываются по критерию \texttt{len(envelope) < 8});
  \item валидную оценку получили 19 сегментов из 79 (24\%).
\end{itemize}

После эксперимента конфигурация возвращена к исходным значениям (25~мс);
изменение \textbf{не было зафиксировано как постоянное}, это разведочный
эксперимент.

\textbf{Вывод эксперимента (не окончательное решение, а констатация факта
для обсуждения с научным руководителем):} оценки $H$ и $T$, получаемые
текущей реализацией, \textbf{нельзя напрямую сравнивать между сегментами
разной длины}, и выбор окна огибающей существенно влияет на результат. Это
не ошибка вычислений — качество аппроксимации $R^2 \geq 0.9$ в обоих
вариантах, — а методическое ограничение текущей конфигурации пайплайна.
Требуется решение: либо нормировать диапазон масштабов $\tau$ независимо от
длины сегмента, либо зафиксировать единое окно огибающей и минимальную
длину анализируемого сегмента по образцу статьи, либо явно учитывать длину
сегмента как ковариату при последующем статистическом анализе.

\textbf{Erratum (31.10.2026).} Решение принято: вариант A+C, см.\
\texttt{docs/rs-methodology-gap-2026-10-31.md}. Реализованы два раздельных
протокола в \texttt{dynamics.py} — \texttt{modified\_rs\_hurst\_full\_record}
(вариант A: огибающая всей записи с паузами, окно 100~мс, геометрическая
шкала $\tau$ до $N$, сопоставим с таблицей~\ref{tab:kuzenkov-results}) и
\texttt{modified\_rs\_hurst} (вариант C: прежний посегментный расчёт,
результат помечен полями \texttt{scope="segment"} и
\texttt{comparable\_to\_reference=False}). Геометрическая шкала $\tau$,
которая в этой редакции отчёта была заявлена, но не реализована, теперь
реализована для обоих протоколов (\texttt{tau\_scale} параметр).

% =====================================================================
\section{Обзор дополнительной литературы}
% =====================================================================

Помимо основной методической статьи~\cite{kuzenkov2014}, для проекта
изучены четыре дополнительных источника, конспекты которых хранятся в
\texttt{docs/books/} и \texttt{docs/literature-review-neuroinformatics-2021.md}:

\begin{itemize}
  \item \textbf{Нейроинформатика-2021} — методы оконной детекции событий
   (YOLOv4+DeepSORT, IoU~78.81\%, F1~0.98), упругие карты для валидации на
   малых выборках через суррогатные данные, спайковые нейросети. Применимо
   к проекту как источник IoU-метрик для оценки временных сегментов и идеи
   робастности через суррогатные данные;
  \item \textbf{HAR на видео (AIP 2025)} — типовой конвейер
   сегментация$\to$признаки$\to$классификация, полу-контролируемое обучение
   при малой разметке, разделение baseline-признаков и обучаемых
   представлений;
  \item \textbf{Wireless HAR (AIP 2026)} — явное обоснование важности
   subject-level (не frame-level) разбиения на train/val/test и deployment-
   метрик (латентность, размер модели);
  \item \textbf{Emotion recognition на DEAP (AIP 2026)} — спектральные
   признаки (FFT, band power, spectral entropy) и табличные baseline-модели
   (XGBoost/RF, accuracy 87--90\%), также подтверждает необходимость
   subject-independent разбиения.
\end{itemize}

Эти источники не являются основой метода (ей остаётся
\cite{kuzenkov2014}), но обосновывают инженерные решения, намеченные на
следующую итерацию: subject-level split, IoU/F1 метрики с tolerance,
табличный baseline поверх классических + нелинейных признаков.

% =====================================================================
\section{Критические ограничения и блокеры}
\label{sec:blockers}
% =====================================================================

\subsection{Блокер 1: отсутствие реального речевого корпуса}

Все результаты раздела~\ref{sec:experiment} получены на синтетических
данных и на \emph{одной} тестовой записи без подтверждённой клинической или
когнитивной разметки. Текущая частота дискретизации 16\,кГц, заложенная в
конфигурации, помечена в коде как допущение
(\texttt{ASSUMPTION: sample\_rate 16000}) и не подтверждена фактическими
данными корпуса. Без реального корпуса невозможны: калибровка порогов VAD,
оценка метрик качества, проверка целевой гипотезы о связи $H/T$ с
интеллектуальными способностями.

\subsection{Блокер 2: отсутствие золотой (gold) разметки}

Нет ручной разметки речевых сегментов на реальных данных, соответственно
невозможен расчёт precision/recall/F1, невозможна калибровка порогов
энергетического VAD и сравнение с baseline-методами на содержательных
данных.

\subsection{Блокер 3: отсутствие baseline WebRTC VAD}

Интерфейс \texttt{VadBackend} подготовлен, но промышленный baseline
(WebRTC VAD, опционально Silero VAD) не реализован — на данный момент
невозможно сравнение собственного подхода с общепринятым решением.

\subsection{Блокер 4: отсутствие разбиения по респондентам (subject-level
split)}

Алгоритм разбиения train/validation/test по идентификатору респондента
(а не по кадрам/сегментам) не реализован. Без него любая оценка качества
на реальных данных будет методологически некорректной (оптимистично
смещённой из-за утечки информации между выборками).

\subsection{Методическое ограничение, выявленное экспериментально}

Зависимость оценок $H$ и $T$ от длины сегмента и от окна огибающей
(раздел~\ref{sec:experiment}) требует явного решения перед тем, как эти
признаки можно будет использовать для сравнения сегментов разной
длительности или для корреляции с когнитивными тестами.

% =====================================================================
\section{Вопросы к научному руководителю}
\label{sec:questions}
% =====================================================================

Для продолжения работы требуются решения по следующим пунктам (полный
список — 12 пунктов в \texttt{docs/questions-for-advisor.md}; ниже —
сокращённая версия для обсуждения).

\subsection{Критически важные}

\begin{enumerate}
  \item \textbf{Реальный корпус.} Когда и как можно получить доступ к
   реальному корпусу записей? Нужны: минимум 50--100 файлов для пилота
   (200+ для валидации), параметры записи (частота дискретизации, число
   каналов, битность, формат), метаданные (\texttt{subject\_id},
   \texttt{task\_type} — чтение/монолог/ответы на вопросы, возраст, пол,
   результаты когнитивных тестов).
  \item \textbf{Разметка.} Согласна ли предложенная таксономия
   (\texttt{speech}, \texttt{silent\_pause}, \texttt{hesitation},
   \texttt{artifact}; дыхание~$\to$~\texttt{artifact}; лексические
   заполнители не выделяются отдельно)? Кто выполняет пилотную двойную
   разметку, сколько файлов, какой целевой уровень согласованности
   (предлагается Cohen's $\kappa \geq 0.8$)?
  \item \textbf{Правовой статус.} Есть ли этическое одобрение на
   использование записей? Разрешено ли использовать внешние облачные
   сервисы ASR/VAD (сейчас действует ограничение
   \texttt{RAW\_AUDIO\_EXTERNAL\_UPLOAD~=~forbidden\_until\_legal\_review})?
   Label Studio — локально или в облаке?
\end{enumerate}

\subsection{Методические}

\begin{enumerate}
  \item \textbf{Окно огибающей и диапазон масштабов $\tau$.} Какое окно
   считать корректным — 25~мс (текущее по умолчанию) или 100~мс (как в
   \cite{kuzenkov2014})? Как бороться с зависимостью $H/T$ от длины
   сегмента (раздел~\ref{sec:experiment}): нормировать диапазон $\tau$,
   фиксировать минимальную длину сегмента, или учитывать длину как
   ковариату?
  \item \textbf{Характерное время $T$.} Достаточен ли реализованный
   кусочно-\emph{линейный} fit (две независимые прямые с оптимальной точкой
   излома), или необходима кусочно-\emph{полиномиальная} аппроксимация
   (линейный слева, квадратичный справа), как указано в оригинальной
   статье?
  \item \textbf{Приоритет дополнительных нелинейных признаков.} Достаточно
   ли на ближайшем этапе довести modified R/S до полной методики
   (включая сравнение с DFA-оценкой Хёрста), или необходимо сразу вводить
   показатель Ляпунова, выборочную/перестановочную энтропию и RQA?
  \item \textbf{Первичная метрика качества VAD.} Оставить ведущей
   \texttt{false\_pause\_precision}, или пересмотреть набор (контрольные:
   \texttt{speech\_recall}, \texttt{hesitation\_recall},
   \texttt{pause\_duration\_bias}; дополнительные: IoU, latency, segment
   F1 с заданным tolerance)?
  \item \textbf{Связь с интеллектуальными способностями.} Какие когнитивные
   тесты планируется использовать для проверки связи (IQ-тесты, тесты
   рабочей памяти, исполнительных функций, беглости речи)? На каком этапе
   — после стабилизации VAD или параллельно? Какие конфаундеры
   контролировать (возраст, образование, нативность языка)?
\end{enumerate}

\subsection{Технические}

\begin{enumerate}
  \item Нужен ли WebRTC VAD как обязательный baseline, или достаточно
   сравнения с energy VAD на первом этапе?
  \item Нужен ли Silero VAD (нейросетевой baseline, требует PyTorch) хотя бы
   как опциональный?
  \item Какие форматы аудио ожидать в реальном корпусе (mono/stereo,
   частота дискретизации, WAV/MP3) — определяет необходимость ресемплинга и
   downmix, которые сейчас не реализованы.
\end{enumerate}

% =====================================================================
\section{Заключение}
% =====================================================================

Инженерный фундамент проекта реализован в соответствии с методической
схемой из статьи~\cite{kuzenkov2014}: модифицированный R/S-анализ с
геометрической шкалой масштабов, усреднением по перекрывающимся окнам и
кусочно-линейной оценкой характерного времени $T$; энергетический VAD с
гистерезисом; единые контракты данных и временная сетка; синтетический
генератор с известной разметкой для проверки корректности алгоритмов;
15 детерминированных unit-тестов, все проходящие.

Технически инструмент работоспособен и воспроизводим. Однако это
подтверждает только \emph{корректность реализации относительно собственной
спецификации на контролируемых сигналах} — не качество детекции речи и не
релевантность признаков $H/T$ на реальных данных, и тем более не связь с
интеллектуальными способностями, которая пока не проверялась ни на каких
данных.

Дополнительно, собственный вычислительный эксперимент
(раздел~\ref{sec:experiment}) выявил, что текущие оценки $H$ и $T$
систематически зависят от длины анализируемого сегмента и от выбора окна
огибающей — это методическая проблема, требующая решения прежде, чем
признаки можно использовать для содержательного сравнения.

Переход к научно значимым результатам блокирован отсутствием: реального
речевого корпуса, золотой разметки, промышленного VAD-baseline и алгоритма
разбиения по респондентам. Следующая итерация работы должна начаться с
решения этих блокеров совместно с научным руководителем (вопросы в
разделе~\ref{sec:questions}); параллельно, без реальных данных, можно вести
реализацию WebRTC-адаптера, DFA-Hurst сравнения на синтетике и подготовку
инструкции разметки.

% =====================================================================
\begin{thebibliography}{9}

\bibitem{kuzenkov2014}
Кузенков~Н.\,П., Логинов~В.\,М.
\textit{Использование метода нормированного размаха при анализе речевых
патологий неврологического генеза.}
Компьютерные исследования и моделирование, 2014, т.~6, №~5, с.~775--791.

\end{thebibliography}

\end{document}
